\documentclass[11pt]{article}
\usepackage[T1]{fontenc}
\usepackage[utf8]{inputenc}
\usepackage{lmodern,amsmath,amssymb,amsthm,mathtools}
\usepackage[margin=1in]{geometry}
\usepackage{microtype}
\usepackage[colorlinks=true,linkcolor=blue,citecolor=blue,urlcolor=blue]{hyperref}
\hypersetup{pdftitle={Failures of unbounded Bayesian prediction: a tanh network and Gaussian-mixture priors},pdfauthor={ChatGPT-6 Astra; Claude Opus 5.5; Kaarel H\"anni; Dmitry Vaintrob}}
\newtheorem{theorem}{Theorem}[section]
\newtheorem{lemma}[theorem]{Lemma}
\newtheorem{corollary}[theorem]{Corollary}
\newcommand{\KL}{\operatorname{KL}}
\title{Failures of unbounded Bayesian prediction:\\a tanh network and Gaussian-mixture priors}
\author{ChatGPT-6 Astra \and Claude Opus 5.5 \and Kaarel H\"anni \and Dmitry Vaintrob\thanks{Authors are listed in alphabetical order.}}
\date{18 November 2026}
\begin{document}
\maketitle
\begin{abstract}
We give an explicit obstruction to distribution-uniform approximation
oracles for the raw posterior mean of a Bayesian neural network.
The model consists of one biased tanh neuron, with three independent
Student weights of any fixed tail index $p>2$, and an ordinary Gaussian
likelihood of any fixed positive variance. The prior has finite second
moments and is unchanged with sample size. There are noiseless
distributions on $[-1,1]\times[-1,1]$ for which the posterior mean has
test squared risk at least $c n^{1/3}$ with constant probability,
although the fixed zero network has risk at most $n^{-1/3}$.
The same divergence holds in expectation. The distributions depend on
$n$ and put mass $1/n$ on a rarely observed test input. The proof
controls the posterior normalizer in both the nearly saturated and
locally linear regimes, and bounds the signed prediction numerator;
it is a lower bound for the actual posterior mean, rather than for
posterior-averaged risk. Thus bounded labels and finite-variance
heavy-tailed coordinates alone do not justify removing an output link
from a random-design Bayesian oracle theorem.
The same construction works with unit-norm inputs in dimension two,
including smooth atomless input laws, and a uniform refinement covers
polylogarithmically growing deterministic likelihood variances.

A separate generic example shows that even a uniformly sub-Gaussian
function supremum and bounded conditional Gaussian covariance do not
suffice for a fixed-temperature agnostic raw-mean oracle. One fixed
mixture of centered rank-one Gaussian function priors on a compact
interval fails at every fixed temperature: its risk grows linearly
along a sample-size sequence, while the zero comparison has fixed
risk and constant relative entropy. The adverse laws may depend on
the temperature. An extension permits arbitrarily light unbounded
amplitude tails, including arsinh-Gaussian amplitudes, with potentially
slower divergence. These generic examples are not finite neural-network
priors. All results concern varying adverse data laws, not inconsistency
under one fixed law.
\end{abstract}

\section{Introduction}
Heavy-tailed weight priors can assign substantial probability to small
neural networks with large coefficients. This can be useful for
approximation, but an unbounded readout creates a separate prediction
problem. A function may fit observed inputs well and have a large
value at an input that the sample misses. A Gaussian likelihood only
penalizes its values on the observed inputs.

We show this problem within a standard finite neural model. There is
one input coordinate, one tanh hidden neuron with a bias, and one
linear readout coefficient. All three parameters are independent
Student variables with tail index $p>2$. Thus the prior is continuous,
has full parameter support and finite second moments, and uses no
architecture mixture. The learner is its ordinary posterior mean
under the Gaussian likelihood using all observations, with any fixed
positive observation variance. Both likelihood and prediction use
the raw network output.

For every sufficiently large $n$, we construct a bounded, noiseless
distribution $P_n$ such that
\[
 P_n^{\otimes n}\{R_{P_n}(\bar f_{S_n})\ge c n^{1/3}\}\ge q,
 \qquad R_{P_n}(0)\le n^{-1/3},
\]
where $c,q>0$ depend only on the fixed prior and likelihood variance.
Here $R_P(h)=\mathbb E_P[(h(X)-Y)^2]$ is population risk at a fresh
random input. The comparator is the same zero function for every $n$.
The bound therefore rules out a distribution-uniform vanishing
oracle remainder, even against this single elementary comparator.
It also rules out such a theorem in expectation.

The construction has two common inputs very close to zero, with
small opposite labels, and a third input of probability $1/n$ with
label zero. With constant probability the sample misses the third
input and contains balanced numbers of the first two. Fitting those
two labels requires a very large local slope. That slope can be
split between the incoming and outgoing weights. The two matching
Student tails and the likelihood Jacobian leave a range of readout
amplitudes whose contribution to the posterior mean is logarithmic
in the required slope. We choose that logarithm to be of order
$n^{2/3}$. The rare input then contributes order $n^{1/3}$ to test
risk. A separate integration argument controls neurons that fit by
saturation; this is needed to upper-bound the full normalizer.

The scope is deliberately precise. The distributions vary with $n$;
we do not establish inconsistency under one fixed distribution.
Their common inputs approach one another rapidly, and we impose no
fixed input regularity constant. The theorem concerns this
width-one Student prior, rather than every heavy-tailed prior or
every larger architecture. Its conclusion is that finite coordinate
variance and bounded labels alone are insufficient for the proposed
raw-mean oracle. Output links, output-norm restrictions, or additional
input and prior structure can materially change the problem.

Appendix~\ref{sec:gaussian-mixture-obstruction} isolates a different
limitation of generic raw-prediction arguments. Its prior is fixed on
smooth functions on $[0,1]$, has a Gaussian tail for the supremum, and
is a mixture of centered Gaussian processes with covariance diagonal
at most one. Nevertheless its raw posterior mean has divergent risk
against a fixed-risk, constant-entropy zero comparator. The latent
features can suppress their values at observed zero-label inputs
while retaining a positive average at unobserved inputs. This is a
general function-prior construction, not another standard neural
architecture. It shows why a fixed-kernel Gaussian argument cannot
be extended using only conditional covariance bounds and native
output tails.

Appendix~\ref{sec:general-amplitude-obstruction} extends this generic
obstruction to every symmetric full-support amplitude law with finite
first moment. It retains that law's native supremum-tail bounds, even
when its moments grow only logarithmically. The failure is still at
fixed temperature, so it is compatible with positive raw-prediction
theorems whose temperature grows as the square of a logarithm.

\paragraph{Related work.}
Dalalyan and Tsybakov~\cite{dt} prove sharp PAC--Bayesian aggregation
bounds for squared loss at a deterministic design. Their later
analysis~\cite{dt-sparse} explicitly accommodates unbounded regression
functions and heavy-tailed priors. Those guarantees concern the
average error on the specified design; the present obstruction
concerns population error at a fresh random input that can be absent
from the sample. There is no conflict between the two statements.
Castillo and Egels~\cite{ce} establish contraction results for dense
deterministic neural architectures with independent heavy-tailed
weights. Such positive results under their statistical assumptions
do not assert the distribution-uniform agnostic raw-mean oracle
refuted here. Our contribution is the explicit fixed neural model
and the direct control of its signed posterior mean; we make no
general claim that negative results for Bayesian prediction are new.

% Independent raw-output obstruction module, 18 November 2026.
% Construction suggested by root; saturation estimate independently checked
% by fresh_reader_approximation and robust_agnostic.
\section{A raw posterior mean can fail under finite-variance Student weights}

Fix $p>2$ and $T>0$. Let the three parameters $v,w,\gamma$ be
independent with density
\[
 q_p(t)=c_p(1+t^2)^{-(p+1)/2},\qquad t\in\mathbb R,
\]
and consider the ordinary biased one-neuron network
$f_{v,w,\gamma}(x)=v\tanh(wx+\gamma)$.
Every coordinate has a finite second moment. The prior, architecture,
and likelihood temperature are fixed independently of the sample size.
For observations $S_n$, use the raw Gaussian likelihood and raw mean
\[
 dQ_{S_n}(\theta)=Z_{S_n}^{-1}
  \exp\left\{-\frac1T\sum_{i=1}^n(f_\theta(X_i)-Y_i)^2\right\}
  d\pi(\theta),\qquad
 \bar f_{S_n}(x)=\int f_\theta(x)\,dQ_{S_n}(\theta).
\]
The physical observation variance is $T/2$. There is no clipping.

\begin{theorem}[A distribution-uniform raw-mean obstruction]
There are constants $c_*,q_*>0$, depending only on $p,T$, and
distributions $P_n$ supported on $[-1,1]\times[-1,1]$, with noiseless
labels, such that for every sufficiently large $n$,
\[
 P_n^{\otimes n}\{R_{P_n}(\bar f_{S_n})\ge c_*n^{1/3}\}\ge q_*.
\]
Yet the fixed width-one comparator $g\equiv0$ satisfies
$R_{P_n}(g)\le n^{-1/3}$. In particular,
$\mathbb E_{P_n^{\otimes n}}R_{P_n}(\bar f_{S_n})\ge q_*c_*n^{1/3}$.
Thus this fixed prior and fixed likelihood do not admit any
distribution-uniform vanishing oracle remainder against the zero
network, either in expectation or at any fixed failure probability
$\delta<q_*$.
\end{theorem}

\begin{proof}
Write
\[
 b=n^{-1/6},\quad c=\frac1{100p(T+1)},\quad
 R=e^{cn^{2/3}},\quad a=b/R.
\]
The distribution assigns probabilities $(1-1/n)/2,(1-1/n)/2,1/n$
to the pairs $(a,b),(-a,-b),(1,0)$, respectively.
Let $E_n$ be the event that the sample contains no input $1$ and that
its counts $n_+,n_-$ at $a,-a$ both belong to $[n/3,2n/3]$.
Its probability tends to $e^{-1}$, so is bounded below by a positive
constant. All bounds below are uniform over samples in $E_n$.

For such a sample define
\[
 u_+=v\tanh(\gamma+aw),\quad
 u_-=v\tanh(\gamma-aw),\quad
 L=\exp\{-[n_+(u_+-b)^2+n_-(u_-+b)^2]/T\}.
\]
The transformation $(v,w,\gamma)\mapsto(-v,-w,-\gamma)$ leaves
the represented function unchanged and preserves the prior. We may
therefore restrict all integrals to $v>0$ and multiply by two.
Constants in the proof may depend on $p,T$ but never on $n$ or the
balanced counts.

Set
\[
 G=\{|u_+-b|\le b/2,\ |u_-+b|\le b/2\},\qquad
 Z_*:=\frac{a^p b^{-p-1}}n=\frac{R^{-p}}{bn}.
\]
On $G^c$, the likelihood is at most
\begin{equation}\label{eq:raw-student-bad-fit}
 L\le \exp\{-nb^2/(12T)\}=\exp\{-n^{2/3}/(12T)\}.
\end{equation}
Since $pc=1/[100(T+1)]<1/(12T)$, this upper bound is
$o(Z_*)$, and also $o(Z_*\log R)$.

We first prove the evidence bound $Z_{S_n}\le C Z_*$.
On $G$ necessarily $v\ge b/2$. For $b/2\le v\le3b$, set
$h_+=\tanh(\gamma+aw)$ and $h_-=-\tanh(\gamma-aw)$.
If the good-fit set of such $v$ is nonempty, then
$h_+,h_-\ge1/6$, whence $w\ge t_0/a$ for
$t_0=\operatorname{arctanh}(1/6)>0$.
For every fixed such $(w,\gamma)$, the exponent as a function of $v$
is a quadratic with coefficient
\[
 n_+h_+^2+n_-h_-^2\ge n/36.
\]
The density $q_p(v)$ is bounded, so integrating in $v$ first, even
over the entire real line, gives at most $C/\sqrt n$.
The Student tail bound then gives
\begin{equation}\label{eq:raw-student-saturation}
 \int_{G\cap\{b/2\le v\le3b\}}L\,d\pi
 \le\frac C{\sqrt n}\int_{t_0/a}^\infty q_p(w)\,dw
 \le \frac{C a^p}{\sqrt n}.
\end{equation}
The bias density integrates to one. This argument includes the
near-saturation region without a singular change of variables.
The ratio of the last bound to $Z_*$ is
$b^{p+1}\sqrt n=n^{(2-p)/6}\le1$ for large $n$.

For $v\ge3b$, use the exact coordinates $(u_+,u_-)$ instead of
$(w,\gamma)$:
\[
 w=\frac{\operatorname{arctanh}(u_+/v)
             -\operatorname{arctanh}(u_-/v)}{2a},\qquad
 \gamma=\frac{\operatorname{arctanh}(u_+/v)
             +\operatorname{arctanh}(u_-/v)}2.
\]
Their Jacobian is
\begin{equation}\label{eq:raw-student-jacobian}
 \left|\frac{\partial(w,\gamma)}{\partial(u_+,u_-)}\right|
 =\frac1{2av^2[1-(u_+/v)^2][1-(u_-/v)^2]}.
\end{equation}
On $G$ we have $|u_\pm|/v\le1/2$ and
$R/(2v)\le w\le2R/v$. Thus the Jacobian is at most $C/(av^2)$.
The global bound $q_p(w)\le C|w|^{-p-1}$ and boundedness of
$q_p(\gamma)$ imply
\[
 q_p(v)q_p(w)q_p(\gamma)
 \left|\frac{\partial(w,\gamma)}{\partial(u_+,u_-)}\right|
 \le C a^p b^{-p-1}q_p(v)v^{p-1}.
\]
The $v$ integral of $q_p(v)v^{p-1}$ over $(0,\infty)$ is finite.
The likelihood integral over the $u$ rectangle is at most
$\pi T/\sqrt{n_+n_-}\le C/n$.
Together with \eqref{eq:raw-student-bad-fit} and
\eqref{eq:raw-student-saturation}, this proves
\begin{equation}\label{eq:raw-student-evidence}
 Z_{S_n}\le C Z_*.
\end{equation}

We next lower-bound the numerator of the raw prediction at $x=1$.
Let $t_n=1/(10\sqrt n)$ and restrict the transformed integral to
\[
 1\le v\le R/100,\quad |u_+-b|\le t_n,\quad
 |u_-+b|\le t_n.
\]
For large $n$ this region lies inside $G$; its likelihood is at least
$e^{-1/(100T)}$. Here $R/(2v)\le w\le2R/v$, $|\gamma|\le1$,
and hence $f(1)=v\tanh(w+\gamma)\ge c_1v$.
The Student density lower bounds give
\[
 q_p(v)\ge c_2v^{-p-1},\quad
 q_p(w)\ge c_2(R/v)^{-p-1},\quad q_p(\gamma)\ge c_2.
\]
The Jacobian in \eqref{eq:raw-student-jacobian} is at least
$1/(2av^2)$. The prediction-weighted density is consequently at
least $c_3 a^p b^{-p-1}/v$ throughout this region. Its $u$ area is
$4t_n^2=1/(25n)$, and its $v$ integral contains
$\int_1^{R/100}dv/v=\log(R/100)$. Its positive contribution to the
numerator is therefore at least $c_4 Z_*\log R$.

There is no negative contribution on $G$. Indeed, for $v>0$ the
opposite signs of $u_+,u_-$ imply $w>0$ and $-aw<\gamma<aw$.
Since $a<1$, it follows that $w+\gamma>w(1-a)>0$.
On $G^c$, $|f(1)|\le|v|$ and
\eqref{eq:raw-student-bad-fit} bounds the absolute contribution by
$\mathbb E_\pi|v|e^{-n^{2/3}/(12T)}=o(Z_*\log R)$.
After symmetry is restored, the full numerator is still at least
$c_5Z_*\log R$. By \eqref{eq:raw-student-evidence},
\[
 \bar f_{S_n}(1)\ge c_6\log R=c_6c n^{2/3}
 \qquad(S_n\in E_n).
\]
The rare input alone contributes at least $(c_6c)^2n^{1/3}$ to test risk.
The zero comparator has risk $(1-1/n)b^2\le n^{-1/3}$, completing
the proof. All posterior quantities exist: the likelihood is bounded
by one, its normalizer is positive, and $|f(x)|\le|v|$ with
$\mathbb E_\pi v^2<\infty$.
\end{proof}

\begin{corollary}[A cheap comparison law still does not give a raw oracle]
Let $\rho_n=\pi(\cdot\mid |v|\le1/n)$. For every data distribution
in the theorem,
\[
 \operatorname{KL}(\rho_n\|\pi)\le\log n+C_p,
 \qquad
 \int R_{P_n}(f)\,d\rho_n(f)
       \le R_{P_n}(0)+n^{-2}.
\]
Consequently no fixed constant $A<\infty$ can make the naive raw-mean
PAC--Bayesian bound
\[
 \mathbb E R_{P_n}(\bar f_{S_n})
 \le\inf_{\rho\ll\pi}
 \left\{\int R_{P_n}(f)\,d\rho(f)
                 +\frac{AT}{n}\operatorname{KL}(\rho\|\pi)\right\}
\]
valid for all bounded data distributions and all $n$ for this prior
and fixed temperature.
\end{corollary}
\begin{proof}
The density $q_p$ is bounded below on $[-1,1]$, so
$\pi(|v|\le1/n)\ge c_p'/n$, giving the entropy bound.
Conditioning on this event preserves the symmetry and independence
of the readout $v$. Thus $\int f(x)d\rho_n=0$ and
$\int f(x)^2d\rho_n\le n^{-2}$ at every input.
Expanding squared loss gives the risk bound. The displayed
PAC--Bayesian right side would then tend to zero, whereas the main
theorem makes its left side diverge. The comparison measure is
deterministic and uses no training data.
\end{proof}

\paragraph{Scope.}
The adverse distribution depends on $n$. This is a failure of a
distribution-uniform oracle theorem, not a claim of inconsistency for
one fixed distribution. The common inputs approach each other
exponentially fast, so the construction imposes no fixed geometric
regularity or uniform anti-concentration constant on the input law.
The teacher is the same zero function for every $n$. The prior is
dense, continuous, independent across all three coordinates, and
unchanged with $n$; the likelihood is the ordinary all-data Gaussian
likelihood with fixed variance. The lower bound is specifically for
the raw posterior mean, rather than for posterior-averaged loss.

\section{The same raw-mean obstruction on the unit circle}
\label{sec:raw-student-circle}
The preceding example does not rely on different input norms. A fixed two-dimensional sphere already permits the same divergence.

\begin{corollary}[A unit-norm input counterexample]
\label{cor:raw-student-circle}
The raw-mean obstruction holds for the one-neuron network
\[
 f_{v,w_1,w_2,\gamma}(x)=v\tanh(w_1x_1+w_2x_2+\gamma),
\]
with four independent coordinates of density $q_p$, when every input lies on the unit circle. The prior, architecture and Gaussian likelihood variance remain fixed. There are noiseless bounded-label laws for which the raw posterior mean has squared test risk at least $c n^{1/3}$ with fixed positive probability, while the zero comparator has risk at most $n^{-1/3}$.
\end{corollary}
\begin{proof}
Retain $b=n^{-1/6}$, $R=e^{c n^{2/3}}$, and $a=b/R$ from the scalar construction, with its sufficiently small fixed $c$. Put $r=\sqrt{1-a^2}$. Replace the common inputs by
\[
 x_+=(a,r),\qquad x_-=(-a,r),
\]
with labels $b,-b$ and probabilities $(1-1/n)/2$ each. The rare input is $x_*=(1,0)$, with label zero and probability $1/n$. Every input has norm one. For large $n$, $r\in[1/2,1]$.

On the no-rare, balanced-count event, the training likelihood depends on $w_2,\gamma$ only through $h=rw_2+\gamma$. The joint density of $(h,\gamma)$ is
\[
 j_r(h,\gamma)=q_p(\gamma)q_p((h-\gamma)/r)/r.
\]
Its marginal $q_r(h)=\int j_r(h,\gamma)d\gamma$ is a probability density, is bounded above uniformly in $r\in[1/2,1]$, and is bounded below by a positive constant on $|h|\le1$, uniformly in that range. The upper bound follows from convolution with a bounded density. For the lower bound, integrate over $|\gamma|\le1$ and use the positive minima of $q_p$ on fixed compact intervals.

Consequently the evidence upper bound from the scalar construction applies unchanged, with $h$ replacing the hidden bias: the saturation estimate uses only that its density integrates to one, and the linear-region estimate uses only a uniform density upper bound. Thus
\[
 Z_{S_n}\le C Z_*,\qquad Z_*=R^{-p}/(bn).
\]
The same global sign symmetry permits restriction to $v>0$.

It remains to check the signed numerator at $x_*$. Define
\[
 I_r(w,h)=\int\tanh(w+\gamma)j_r(h,\gamma)d\gamma.
\]
On the good-fit rectangle, $w=w_1>0$ and $|h|<aw$, so $w+h>0$. Pair the integral around $\gamma=-w$. For $t>0$, monotonicity of $q_p$ in absolute value gives
\begin{align*}
 q_p(-w+t)&\ge q_p(-w-t),\\
 q_p((h+w-t)/r)&\ge q_p((h+w+t)/r).
\end{align*}
Both inequalities use positive centers, $w>0$ and $h+w>0$. Therefore
\[
 I_r(w,h)=\int_0^\infty\tanh(t)
       [j_r(h,-w+t)-j_r(h,-w-t)]dt\ge0.
\]
Thus the full good-fit region has no negative contribution after integrating the latent bias decomposition.

Moreover, there are fixed $W_0,c_0>0$ such that
\[
 I_r(w,h)\ge c_0\qquad(w\ge W_0,\ |h|\le1,\ 1/2\le r\le1).
\]
To verify this, use only the paired terms with $\gamma=-w+t\in[-1,1]$. Their positive density $j_r(h,\gamma)$ is bounded below by an absolute positive constant. The reflected density at $-2w-\gamma$ is at most $Cw^{-2p-2}$, uniformly in $h,r$. For sufficiently large fixed $W_0$ it is at most half the positive lower bound, while $\tanh(t)\ge1/2$. Integration over this interval proves the claim; all omitted paired terms are nonnegative.

In the scalar lower-numerator region, replace $1\le v\le R/100$ by
$1\le v\le R/(100W_0)$. Then $w\ge W_0$, $|h|\le1$, and the same Jacobian and Student bounds apply. The integrated prediction factor is now $v I_r(w,h)\ge c_0v$. Hence the positive numerator is at least $cZ_*\log R$, since shortening the upper endpoint by a fixed factor changes the logarithm only by a constant. The bad-fit contribution is still at most
$\mathbb E|v|e^{-n^{2/3}/(12T)}=o(Z_*\log R)$ in absolute value.
Therefore the raw posterior mean at $x_*$ is at least $c\log R$, and its rare-input risk is at least $c n^{1/3}$. The event probability and zero-comparator risk are unchanged.
\end{proof}

This unit-norm version still uses input laws whose two common atoms approach one another exponentially fast. It does not satisfy a fixed uniform anti-concentration condition. The conclusion is a distribution-uniform failure, not inconsistency for a single fixed input law.

\begin{corollary}[Smooth atomless input distributions]
The input laws in Corollary~\ref{cor:raw-student-circle} can instead
have $C^\infty$ densities with respect to arc length on the unit
circle. Labels remain noiseless and bounded. The same risk lower
bound, up to a fixed constant factor, and the same zero-comparator
risk upper bound hold. No uniform bound on the input densities or
their derivatives is asserted.
\end{corollary}
\begin{proof}
Fix $n$ and its three-point construction. For every parameter vector,
$f_\theta(x)$ is continuous in $x$ on the circle. The posterior
normalizer is continuous in the training inputs and labels, by
dominated convergence with the integrable bound one. The prediction
numerator is jointly continuous in the training data and test input,
by dominated convergence with the integrable bound $|v|$.
The normalizer is strictly positive everywhere. On the compact set
of circle-valued training inputs, labels in $[-1,1]$, and a
circle-valued test input, it therefore has a positive minimum.
Consequently the raw posterior prediction is jointly continuous
and uniformly continuous on this compact set. These statements are
for each fixed $n$ and its positive temperature; no uniform modulus
in $n$ is needed.

Replace each of the three input atoms by a normalized smooth bump
supported in a sufficiently small arc around that atom, retaining
its original mixture probability. Choose the three arcs disjoint,
and retain the constant label $b,-b,0$ on the respective arcs.
Couple each perturbed observation with its original cluster center.
For sufficiently small arc radii, uniform continuity ensures that
the perturbed posterior prediction, at every test input in the rare
arc, differs from the original posterior prediction at $x_*$ by
at most one, simultaneously for every training cluster assignment.
On the original no-rare, balanced-count event, the original
prediction is at least $c n^{2/3}$. The perturbed prediction is
therefore at least $(c/2)n^{2/3}$ for all sufficiently large $n$.
The event probability is unchanged because it depends only on the
cluster assignments. The rare arc still has probability $1/n$ and
label zero, so it contributes at least $(c^2/4)n^{1/3}$ to risk.
The zero-comparator risk remains exactly $(1-1/n)b^2$.
The finite mixture of the three smooth bump densities is smooth
and atomless, proving the claim.
\end{proof}

The arc radii may depend arbitrarily strongly on $n$. This smoothing
does not give a counterexample under a fixed density bound or a
fixed quantitative anti-concentration condition.

\section{Deterministic growing likelihood variances}
\label{sec:raw-growing-temperature}
The obstruction persists for slowly growing, predetermined likelihood
variances. The following statement allows a different adverse law at
each sample size, as does the main theorem. The temperature sequence
must be fixed before observing the data.

\begin{corollary}[A growing-temperature obstruction]
Fix $p>2$, $T_0>0$, and
\[
 \frac12<r<\frac p{p+1}.
\]
Let $(T_n)$ be any deterministic sequence satisfying
$T_0\le T_n\le n^{r-1/2}$ for all sufficiently large $n$.
For the same fixed one-neuron Student prior and the raw Gaussian
likelihood at temperature $T_n$, there are bounded noiseless laws
$P_n$ and constants $c_*,q_*>0$, depending only on $p,r,T_0$, such that
\[
 \begin{split}
 P_n^{\otimes n}\left\{R_{P_n}(\bar f_{S_n})
       \ge c_*\frac{n^{2r-1}}{(T_n+1)^2}\right\}&\ge q_*,\\
 R_{P_n}(0)&\le n^{r-1}.
 \end{split}
\]
The same lower bound, multiplied by $q_*$, holds in expectation.
In particular, the risk diverges for every polylogarithmically
bounded temperature sequence, whereas the zero-comparator risk
vanishes. For $p>3$, one may choose $r>3/4$; then the risk still
diverges when $T_n=n^{1/4}$.
\end{corollary}

\begin{proof}
Use the same three-point construction and balanced-count event, but
now set
\[
 b=n^{(r-1)/2},\qquad
 R=\exp\left\{\frac{n^r}{100p(T_n+1)}\right\},\qquad a=b/R.
\]
All prior density bounds in the preceding proof are independent of
temperature. Integrating the likelihood in the saturated region
now gives $C\sqrt{T_n/n}$, so that contribution is bounded by
$C a^p\sqrt{T_n/n}$. Integrating the two regular prediction
coordinates gives $CT_n/n$. The appropriate evidence scale is
\[
 Z_{*,n}:=\frac{T_n}{n}a^p b^{-p-1}
        =\frac{T_nR^{-p}}{bn}.
\]
The ratio of the saturated bound to this scale is at most
\[
 C b^{p+1}\sqrt{n/T_n}
 \le\frac C{\sqrt{T_0}}
       n^{[(r-1)(p+1)+1]/2}=o(1).
\]
On the bad-fit complement the likelihood is at most
$\exp\{-n^r/(12T_n)\}$.
Its exponential rate exceeds $p\log R=n^r/[100(T_n+1)]$ by
at least $c n^r/(T_n+1)$. Since $T_n\le n^{r-1/2}$, this last
quantity is at least $c'\sqrt n$ for large $n$. The polynomial
factors in $Z_{*,n}$ cannot offset this gap. Hence
$Z_{S_n}\le C Z_{*,n}$ uniformly over the allowed temperatures.

For the numerator use prediction-coordinate radius
$t_n=\sqrt{T_n}/(10\sqrt n)$ instead of $1/(10\sqrt n)$.
Its ratio to $b$ is
\[
 \frac{t_n}b=\frac1{10}\sqrt{T_n/n^r}
                 \le\frac1{10}n^{-1/4},
\]
so the rectangle still lies in the good-fit region. Its likelihood
is at least $e^{-1/100}$ and its area is $T_n/(25n)$.
The same density and Jacobian lower bounds, integrated over
$1\le v\le R/100$, therefore give numerator at least
$c Z_{*,n}\log R$. The negative contribution is still bounded by
$\mathbb E|v|e^{-n^r/(12T_n)}=o(Z_{*,n}\log R)$.
Consequently $\bar f_{S_n}(1)\ge c\log R$, which gives the stated
rare-input risk. All constants are uniform in the prescribed
temperature sequence within the displayed range.
\end{proof}

This corollary does not cover data-dependent temperature selection.
It also does not assert failure for every possible deterministic
temperature growth rate.


\section{Consequences and limits}
The adverse sample event has probability tending to $e^{-1}$.
Consequently the lower bound persists at any fixed confidence level
whose allowed failure probability is smaller than some positive
constant. The target is noiseless, so there is no irreducible noise
floor. The construction works for every fixed positive likelihood
variance, with its constants and adverse distributions allowed to
depend on that variance.

The result also distinguishes the posterior mean from a posterior
draw. A large posterior second moment need not force a large mean,
because positive and negative predictions can cancel. Here the
good-fit region has positive prediction at the rare input, while the
absolute contribution from its complement is exponentially smaller
than the positive numerator. This explicit sign argument is what
turns the tail calculation into a lower bound for the raw mean.

The finite second moment of the prior guarantees that every
finite-sample posterior mean and risk in this construction is finite.
The statement is a divergence along a sequence of sample sizes and
data distributions, rather than an undefined posterior moment.
Clipping each prediction to the label range would prevent this
particular amplitude divergence. The theorem does not, by itself,
establish any oracle guarantee for a clipped learner.

Section~\ref{sec:raw-growing-temperature} gives the corresponding
range of exponents and deterministic temperatures. These variants
retain the same prior and architecture. The likelihood variance
changes with sample size only when the temperature sequence is
explicitly allowed to do so.

The generic examples rule out even fixed multiplicative agnostic
factors under its covariance and tail assumptions. It does not
contradict fixed-temperature realizable guarantees that localize the
prior near the true regression function: its cheap zero comparison
has approximation error $1/2$. Nor does it establish failure of any
particular normalized Gaussian-readout neural prior. A positive
theorem for such a prior may use feature structure beyond the
generic assumptions considered here.

\appendix
\section{A bounded-covariance Gaussian mixture can fail under raw prediction}
\label{sec:gaussian-mixture-obstruction}
A fixed Gaussian function prior and a mixture of Gaussian function priors have different raw-output behavior. The following example has a uniformly sub-Gaussian function supremum. Conditional on its latent feature, it is a centered rank-one Gaussian process with kernel diagonal at most one. Nevertheless its ordinary raw posterior mean has divergent population risk at every prescribed fixed likelihood temperature. The example concerns a general function prior; it is not a counterexample for the specific continuous neuron-gate prior constructed elsewhere.

\begin{theorem}[A fixed-prior Gaussian-mixture obstruction]
\label{thm:gaussian-mixture-obstruction}
For every fixed $T>0$, there is one prior $\mu_T$ on continuous real-valued functions on $[0,1]$, unchanged with sample size, with the following properties.
\begin{enumerate}
 \item Conditional on a latent feature, the function is a centered Gaussian process whose covariance is rank one or zero and whose kernel obeys $K(x,x)\le1$ at every input.
 \item If $F(f)=\sup_{x\in[0,1]}|f(x)|$, then
 $\mathbb E_{\mu_T}\exp(F(f)^2/4)\le\sqrt2$.
 \item The zero function has prior mass $1/2$.
\end{enumerate}
There are noiseless laws $P_j$ with labels in $\{0,1\}$ and sample sizes $n_j=2^j$ such that the ordinary posterior
\[
 \Pi_S(df)\propto
       \exp\left\{-\frac1T\sum_{i=1}^{n_j}(f(X_i)-Y_i)^2\right\}\mu_T(df)
\]
has raw posterior mean $m_S=\int f\,d\Pi_S$ satisfying, for all sufficiently large $j$,
\[
 P_j^{\otimes n_j}\{R_{P_j}(m_S)\ge c_T n_j\}\ge9/10,
 \qquad R_{P_j}(0)=1/2,
\]
where $c_T>0$ depends only on $T$. The expected risk also diverges. All function paths can be chosen smooth, and all finite-sample posterior moments used in the proof exist.
\end{theorem}

\begin{proof}
Fix $T>0$ and put
\[
 p=\frac1{100(T+1)^2},\qquad n_j=2^j,\qquad M_j=n_j^3.
\]
Choose disjoint open intervals $I_j\subset(0,1)$ accumulating only at zero. In each $I_j$, choose $M_j+1$ distinct points $x_{j,0},\ldots,x_{j,M_j}$ and smooth bump functions $\eta_{j,i}$ with pairwise disjoint supports compactly contained in $I_j$, satisfying $0\le\eta_{j,i}\le1$ and $\eta_{j,i}(x_{j,i})=1$.
For a bit vector $B=(B_1,\ldots,B_{M_j})$, define
\[
 \psi_{j,B}(x)=n_j^{-1/2}\eta_{j,0}(x)
                   +\sum_{i=1}^{M_j}B_i\eta_{j,i}(x).
\]
Disjoint supports ensure $0\le\psi_{j,B}\le1$.

Give the zero function prior mass $1/2$. Give branch $j$ mass
$w_j=2^{-j-1}$. Within that branch draw independent bits
$B_i\sim\operatorname{Bernoulli}(p)$ and an independent $G\sim N(0,1)$, and set
$f=G\psi_{j,B}$. These branch masses sum to $1/2$, so this defines one fixed prior. Conditional on $(j,B)$, its covariance is
$K_{j,B}(x,x')=\psi_{j,B}(x)\psi_{j,B}(x')$, with diagonal at most one. Also $F(f)\le|G|$, proving the uniform exponential moment. Each sampled function is smooth and supported in one interval; no path contains an infinite sum of bumps over the branch index.

Under $P_j$, the common input $x_{j,0}$ has probability $1/2$ and label one. Each of the $M_j$ other inputs has probability $1/(2M_j)$ and label zero. Write $n=n_j$, $M=M_j$. Let $K$ be the number of common observations and let $k_i$ be the number of observations at rare input $i$. Define the event
\[
 E_n=\{K\ge n/3\}.
\]
Since $K\sim\operatorname{Binomial}(n,1/2)$, Hoeffding's inequality gives
$\Pr(E_n)\ge1-e^{-n/18}$.

All branches other than $j$ are identically zero on this law's support. Their total likelihood is therefore
$Z_0=e^{-K/T}$. In the active branch, after averaging each observed bit, the conditional likelihood for $G$ is
\[
 e^{-K(G/\sqrt n-1)^2/T}\,L(G),\qquad
 L(G)=\prod_{i:k_i>0}[(1-p)+p e^{-k_iG^2/T}].
\]
There are at most $n$ observed rare coordinates, so
\begin{equation}
 (1-p)^n\le L(G)\le1.
 \label{eq:gp-mixture-bit-tilt}
\end{equation}
Set $k=K/n$. Integrating only the common-observation Gaussian factor gives
\[
 Z_c=(1+2k/T)^{-1/2}\exp\{-K/(T+2k)\}.
\]
Consequently the active-branch to inactive-branch posterior odds are at least
\begin{align*}
 \frac{w_j}{1-w_j}\frac{(1-p)^n Z_c}{Z_0}
 &\ge \frac{w_j}{1-w_j}(1+2/T)^{-1/2}
       \exp\left\{n\left[
         \frac{2k^2}{T(T+2k)}+\log(1-p)\right]\right\}.
\end{align*}
On $E_n$, $1/3\le k\le1$. The bracket is at least
\[
 \frac{2/9-1/50}{(T+1)^2}>0,
\]
because $\log(1-p)\ge-2p$ and
$T(T+2)\le(T+1)^2$. The exponential in $n=2^j$ dominates
$\log(1/w_j)=O(j)$. Hence the active posterior mass is at least $1/2$ for all sufficiently large $j$, uniformly over samples in $E_n$.

We next bound its posterior mean readout. After conditioning only on the common observations, the readout is Gaussian with
\[
 \mu=\frac{2k\sqrt n}{T+2k},\qquad
 v=\frac T{T+2k}\le1.
\]
On $E_n$, $\mu\ge2\sqrt n/[3(T+1)]$ and $\mu\le\sqrt n$.
The active readout posterior $Q$ is this Gaussian tilted by $L(G)$.
By \eqref{eq:gp-mixture-bit-tilt} its density ratio is at most
$(1-p)^{-n}\le e^{2pn}$. Thus
\[
 Q\{G\le\mu/2\}
 \le\exp\left\{-\frac{(1/18-1/50)n}{(T+1)^2}\right\}.
\]
The negative contribution is bounded, by Cauchy--Schwarz under the unperturbed Gaussian, by
\[
 \mathbb E_Q[(-G)_+]
 \le\sqrt{n+1}\,
       \exp\left\{-\frac{(1/9-1/50)n}{(T+1)^2}\right\}.
\]
Indeed the Gaussian probability of $G<0$ is at most
$e^{-\mu^2/(2v)}$, and its second moment is at most $n+1$.
It follows for large $j$ that
\[
 \mathbb E_QG\ge\mu/4\ge\frac{\sqrt n}{6(T+1)}.
\]

Every unobserved bit remains independent Bernoulli$(p)$ under this active posterior and remains independent of $G$: the likelihood contains no occurrence of that bit. Hence at every unobserved rare input the full posterior mean is at least
\[
 \frac12 p\,\mathbb E_QG
       \ge\frac{p\sqrt n}{12(T+1)}.
\]
There are at least $M-n\ge M/2$ such inputs. Their total population probability is at least $1/4$, and their labels are zero. Therefore, on $E_n$,
\[
 R_{P_j}(m_S)\ge\frac{p^2 n}{576(T+1)^2}.
\]
This proves the theorem with $c_T=p^2/[576(T+1)^2]$, since $\Pr(E_n)\to1$. The zero-function risk is exactly $1/2$. The likelihood is bounded by one, its evidence is positive, and the prior function supremum has every finite moment, so all posterior quantities in the proof exist.
\end{proof}

\begin{corollary}[One prior for every fixed temperature]
There is one prior $\mu$, independent of temperature as well as sample
size, with all three properties of
Theorem~\ref{thm:gaussian-mixture-obstruction}, such that its conclusion
holds for every fixed $T>0$. The positive constant and adverse data-law
sequence may depend on $T$; a single data-law sequence defeating all
temperatures is not asserted.
\end{corollary}
\begin{proof}
Use pairwise disjoint input intervals $I_{j,k}$ for all integers
$j,k\ge1$. Give the zero function mass $1/2$ and branch $(j,k)$ mass
$2^{-j-k-1}$; these branch masses sum to $1/2$. On that branch use the
preceding smooth-feature construction with $n_j=2^j$, $M_j=n_j^3$,
and Bernoulli parameter $p_k=2^{-k}/100$. The Gaussian amplitude and
all supremum and covariance bounds are unchanged. For any fixed
$T$, choose one fixed $k=k(T)$ with
$p_k\le1/[100(T+1)^2]$ and support the adverse law on block $(j,k)$.
All other branches vanish on that block, including those with
smaller Bernoulli parameters. Every estimate in the preceding proof
still holds, while the extra prior penalty is only $O(k)$ in
addition to $O(j)$ and is dominated by the exponential evidence
advantage. The final constant is
$c_T=p_{k(T)}^2/[576(T+1)^2]>0$.
\end{proof}

\begin{corollary}[No covariance-only raw agnostic oracle]
There is no distribution-uniform oracle bound for the temperature-independent prior $\mu$ above and any fixed $T>0$ of the form
\[
 \mathbb E R(m_S)\le C_1\inf_{\rho\ll\mu}
       \left\{\rho R+\frac{C_2T}{n}\KL(\rho\|\mu)\right\}+o(1),
\]
with finite constants $C_1,C_2$, or with any other right side that stays bounded for a fixed-risk, fixed-entropy comparator. Indeed the comparison law concentrated on the zero atom has $\rho R=1/2$ and $\KL(\rho\|\mu)=\log2$.
\end{corollary}

The adverse laws vary with $j$, whereas the function prior and likelihood variance are fixed. The theorem does not assert inconsistency under one fixed law. It also does not identify the above feature mixture with a finite neural architecture. Its purpose is to isolate a statistical obstruction: a uniform sub-Gaussian raw supremum and bounded conditional Gaussian covariance do not, by themselves, permit the sharp fixed-kernel Gaussian argument to be averaged over learned feature kernels. The latent feature weights can use observed zero labels to switch off training coordinates while leaving an unobserved positive bias.

\section{Arbitrarily light unbounded tails do not suffice at fixed temperature}
\label{sec:general-amplitude-obstruction}
The preceding Gaussian-mixture construction is not specific to Gaussian
amplitude tails. Even much lighter unbounded tails cannot, by themselves,
imply a general fixed-temperature raw-prediction oracle. The following
is still a generic function-prior example, not a counterexample for a
particular finite neural-network prior.

\begin{theorem}[One fixed prior for every fixed temperature]
\label{thm:general-amplitude-obstruction}
Let $G$ have any symmetric full-support law on $\mathbb R$ with
$\mathbb E|G|<\infty$. There is one centered prior $\mu$ on smooth
functions on $[0,1]$, independent of temperature, such that
\[
 \mathbb E_\mu\Psi(\|f\|_\infty)
       \le\mathbb E\Psi(|G|)
\]
for every nonnegative increasing function $\Psi$.
For every fixed $T>0$ there is a sequence of noiseless data laws,
with labels in $\{0,1\}$ and sample sizes $n_j\to\infty$, for which
the ordinary raw Gaussian-likelihood posterior mean satisfies
\[
 \mathbb P\{R(\bar f_{S_{n_j}})\ge c_Tj^2\}\longrightarrow1.
\]
The zero predictor has risk $1/2$ and a comparison law of entropy
at most $\log2$ under this one prior. In particular the conclusion applies to
$G=\operatorname{arsinh}(Z)$ with $Z$ standard Gaussian, whose moments
grow only logarithmically. The adverse law sequence is allowed to
depend on $T$; one fixed law or one law sequence defeating all
$T$ is not asserted.
\end{theorem}
\begin{proof}
For each integer $j\ge1$, set
$r_j=\mathbb P\{j\le G\le9j/8\}>0$.
Choose an integer
\[
 n_j\ge2^j+j^2\{1+\log(1/r_j)\},\qquad M_j=n_j^3,
\]
increasing with $j$, if necessary by increasing it further.
For each $k\ge1$ let $p_k=2^{-k}/1000$.
Choose disjoint open intervals $I_{j,k}\subset(0,1)$ and, inside
each, $M_j+1$ disjoint nonnegative smooth bump functions
$\zeta_{j,k,i}$ of maximum one. Let $x_{j,k,i}$ be points at which
the respective bumps equal one. For independent Bernoulli$(p_k)$
bits $B_1,\ldots,B_{M_j}$, define
\[
 \psi_{j,k,B}(x)=\frac1j\zeta_{j,k,0}(x)
                         +\sum_{i=1}^{M_j}B_i\zeta_{j,k,i}(x).
\]
The prior has an atom of mass $1/2$ at zero and, for each $(j,k)$,
mass $w_{j,k}=2^{-j-k-1}$ on $f=G\psi_{j,k,B}$, independently
drawing $G$ and the bits. These weights sum to one. Every individual
function is smooth and supported in one interval, and
$\|\psi_{j,k,B}\|_\infty\le1$. Symmetry centers the prior and
proves the asserted tail comparison.

Fix $T>0$ and choose $k$ such that
$p=p_k\le1/[1000(T+1)]$. For each $j$, put probability $1/2$ on
$x_{j,k,0}$ with label one, and probability $1/(2M_j)$ on each
$x_{j,k,i}$, $i\ge1$, with label zero. Every inactive prior branch
is identically zero at these inputs. Let $K$ count the common input
in a sample of size $n=n_j$. The event $K\ge n/3$ has probability
tending to one. We work on this event.

Within the active branch, integrating observed bits gives the
amplitude likelihood
\[
 \exp[-K(G/j-1)^2/T]L(G),\qquad
 L(G)=\prod_{i\,\mathrm{seen}}
       \{(1-p)+p\exp(-k_iG^2/T)\},
\]
where $k_i$ is the observed multiplicity of rare input $i$.
For every $G$, $(1-p)^n\le L(G)\le1$.
Its branch normalizer $D_j$ is therefore bounded below by
\[
 D_j\ge r_j\exp[-n/(64T)](1-p)^n
       \ge r_j\exp[-n/(64T)-2pn].
\]
If $G\le j/2$, the common-input squared error is at least $1/4$.
Consequently the active-branch posterior $Q_j$ satisfies
\[
 Q_j\{G\le j/2\}
 \le r_j^{-1}\exp[-n/(12T)+n/(64T)+2pn]
 \le r_j^{-1}e^{-n/(20T)}.
\]
For negative amplitudes the common-input error is at least one, so
\[
 \mathbb E_{Q_j}[|G|\mathbf1_{\{G<0\}}]
       \le(\mathbb E|G|)r_j^{-1}e^{-n/(4T)}.
\]
The last two numerical bounds follow from
$p\le1/[1000(T+1)]$. Since
$\log(1/r_j)\le n_j/j^2$, both right sides tend to zero.
It follows that $\mathbb E_{Q_j}G\ge j/4$ for all sufficiently
large $j$.

The combined inactive branches have likelihood $e^{-K/T}$.
The posterior odds of the active branch are at least
\[
 \frac{w_{j,k}}{1-w_{j,k}}\frac{D_j}{e^{-K/T}}
       \ge\frac{w_{j,k}}{1-w_{j,k}}r_j e^{n/(4T)}
       \longrightarrow\infty.
\]
Here $-\log w_{j,k}=O(j+k)$, $n_j\ge2^j$, and
$-\log r_j\le n_j/j^2$. Thus active posterior mass is at least
one half eventually. At an unobserved rare input, its bit retains
its independent Bernoulli$(p)$ law, even after conditioning on the
amplitude and all observed data. The full raw posterior mean there is
therefore at least $p j/8$. At least $M_j-n_j$ rare inputs are
unobserved, and their total population mass is at least $1/4$.
The resulting squared risk is at least $p^2j^2/256$, proving the
claim with $c_T=p_k^2/256$.

The zero predictor has exactly risk $1/2$. The zero function has prior probability at least $1/2$, giving the
stated comparison entropy. All posterior
means exist because the likelihood is bounded and
$\mathbb E|G|<\infty$; the displayed lower bounds remain valid if
some population risks are infinite.
\end{proof}

The theorem preserves any prescribed native tail control obtainable
from a symmetric unbounded amplitude law. For the asinh-Gaussian
choice it is therefore compatible with the logarithmic-temperature
positive oracle: it only rules out a generic conclusion at fixed
$T$. Its function prior uses a countable mixture of feature families
for the counterexample; it is not the fixed-architecture neural prior
used in that positive result. The quantitative divergence here is
measured in $j$, and its sample sizes can be much larger than in the
Gaussian-amplitude example.


\begin{thebibliography}{9}
\bibitem{dt} Arnak Dalalyan and Alexandre B. Tsybakov.
Aggregation by exponential weighting, sharp PAC-Bayesian bounds and sparsity.
\emph{Machine Learning} 72(1--2):39--61, 2008.
\url{https://doi.org/10.1007/s10994-008-5051-0}.
\bibitem{dt-sparse} Arnak S. Dalalyan and Alexandre B. Tsybakov.
Sparse regression learning by aggregation and Langevin Monte-Carlo.
\emph{Journal of Computer and System Sciences} 78(5):1423--1443, 2012.
\url{https://doi.org/10.1016/j.jcss.2011.12.023}.
\bibitem{ce} Isma\"el Castillo and Paul Egels.
Posterior and variational inference for deep neural networks with heavy-tailed weights.
\emph{Journal of Machine Learning Research} 26(122):1--58, 2025.
\url{https://www.jmlr.org/papers/v26/24-0894.html}.
\end{thebibliography}
\end{document}
